\chapter{LQR com ação integral (LQI)}
\label{cap:lqi}

\section{Motivação}
O LQR puro tem erro de regime não nulo diante de perturbações constantes: vento, viés
da malha de velocidade da PX4, erro de escala da profundidade, erro na estimativa de
$\vect v_t$. Se a perturbação é uma aceleração constante $\bar a$, o regime do LQR é
$e_{p,ss} = \bar a/k_p$. Um integrador do erro de posição elimina esse resíduo.

\section{Modelo aumentado}
Arquivo \code{control/lqi\_controller.cpp}. Com $\xi_{k+1}=\xi_k+\Ts e_{p,k}$:
\begin{equation}
  z_k=\begin{bmatrix}e_p\\e_v\\ \xi\end{bmatrix}_k,\quad
  A_a=\begin{bmatrix}1&\Ts&0\\0&1&0\\ \Ts&0&1\end{bmatrix},\quad
  B_a=\begin{bmatrix}\Ts^2/2\\ \Ts\\ 0\end{bmatrix},\quad
  Q_a=\mathrm{diag}(q_p,q_v,q_\xi).
\end{equation}
A DARE 3$\times$3 dá $u=-[k_p,\ k_d,\ k_\xi]\,z$. Com o \code{yaml} ($q_\xi=2$):
$K=[5{,}52,\ 3{,}76,\ 1{,}81]$, polos equivalentes $-1{,}77\pm1{,}12j$ e $-0{,}456$. O
polo do integrador (\SI{2,2}{s}) é cerca de 4 vezes mais lento que o par dominante,
que é a separação desejada. O preço é a margem: com PX4 e \SI{0,1}{s} de atraso, a PM
cai de \ang{50} (LQR) para \ang{43} (LQI, contando só $k_p,k_d$; o integrador tira mais alguns graus), com $\omega_c$ subindo de 1,25 para
\SI{1,51}{rad/s}.

\section{Como escolher $q_\xi$}
Com $r=1$ e o par dominante fixado por $q_p,q_v$, o polo do integrador vale
aproximadamente
\begin{equation}
  s_\xi \approx -\frac{k_\xi}{k_p} \approx -\sqrt{\frac{q_\xi}{q_p}}\cdot c,
\end{equation}
em que $c=O(1)$ depende dos demais pesos. Na prática, varra $q_\xi$ no
\code{projeto\_lqi.m} e escolha o valor que põe $s_\xi$ entre 1/3 e 1/5 da parte real
do par dominante. Integrador mais rápido que isso derruba a margem de fase e reaparece
o sobressinal.

\section{Anti-windup}
Durante a aproximação inicial (erro de vários metros, velocidade saturada), o
integrador acumularia um valor enorme e causaria sobressinal ao chegar. O código usa:
\begin{itemize}
  \item \textbf{integração condicional}: só integra se $\lVert\ep\rVert<$
        \code{integral\_zone} (\SI{2}{m});
  \item \textbf{saturação}: $\lVert\xi\rVert\le$ \code{integral\_limit} (\SI{5}{m.s}).
\end{itemize}
Alternativa clássica para o TCC: \emph{back-calculation},
$\dot\xi = e_p + k_{aw}(u_{sat}-u)$, com $k_{aw}\approx 1/T_i$.

\section{Controle ótimo e MATLAB}
O LQI já é ótimo para o custo aumentado. No MATLAB:
\begin{lstlisting}
Aa = [1 Ts 0; 0 1 0; Ts 0 1];  Ba = [Ts^2/2; Ts; 0];
[K,~,polos] = dlqr(Aa, Ba, diag([10 2 2]), 0.5);
% contínuo, com Control System Toolbox:
sys = ss([0 1;0 0],[0;1],[1 0],0);  K = lqi(sys, diag([10 2 2]), 0.5);
\end{lstlisting}
Atenção ao sinal: \code{lqi} integra $r-y$, o oposto de $\xi$ no código; o ganho do
integrador sai com sinal trocado.

\begin{receita}
\textbf{LQI (\code{matlab/projeto\_lqi.m}).} Projete $q_p,q_v,r$ como no LQR. Varra
$q_\xi$ e observe o polo do integrador e a PM com PX4 + atraso. Valide com uma
perturbação constante (vento de \SI{0,5}{m/s^2} no \code{sim\_perseguicao.m}): o LQR fica
com erro $\bar a/k_p$, e o LQI zera o erro em $\approx 4/|s_\xi|$ segundos.
\end{receita}
